\chapter{Reproducibility Checklist}
\label{app:repro}

This appendix provides a reproducibility checklist for replicating our results.

%==============================================================================
\section{Checklist}
%==============================================================================

\begin{table}[h]
\centering
\caption{Reproducibility Checklist}
\begin{tabular}{ll}
\toprule
Item & Status \\
\midrule
Code available & $\checkmark$ (repository) \\
Data source & MIMIC-IV v3.1 (credentialed access) \\
Random seed & 42 (fixed) \\
Hyperparameters & Appendix A (full grids) \\
SQL queries & Appendix B \\
Environment & Python 3.10, PyTorch 2.x, scikit-learn 1.3+ \\
Hardware & CPU sufficient; GPU optional for ETHOS \\
\bottomrule
\end{tabular}
\label{tab:repro_checklist}
\end{table}

%==============================================================================
\section{Code Listings}
%==============================================================================

Representative code for tokenization (\texttt{tokenize.py}):

\begin{lstlisting}[language=Python, caption={Tokenization (excerpt)}]
def tokenize_symptom(value: float, bin_edges: np.ndarray) -> int:
    """Map continuous value to discrete token ID."""
    bin_id = np.searchsorted(bin_edges, value, side='right') - 1
    bin_id = np.clip(bin_id, 0, len(bin_edges) - 2)
    return bin_id + 1  # 0 reserved for mask

def add_intensity(seq: np.ndarray, values: np.ndarray) -> np.ndarray:
    """Append intensity as extra dimension."""
    return np.concatenate([seq, values[:, None]], axis=-1)
\end{lstlisting}

%==============================================================================
\section{Data Availability}
%==============================================================================

MIMIC-IV requires credentialed access via PhysioNet. Cohort extraction scripts are provided. Processed parquet files can be regenerated from raw MIMIC-IV using the pipeline in \texttt{scripts/}.

%==============================================================================
\section{Environment and Dependencies}
%==============================================================================

To reproduce our results, use the following environment:

\begin{lstlisting}
Python 3.10+
torch>=2.0.0
scikit-learn>=1.3.0
pandas>=2.0
numpy>=1.24
xgboost>=2.0
lightgbm>=4.0
catboost>=1.2
\end{lstlisting}

Run experiments from the repository root (requires MIMIC-IV under \texttt{data/3.1/} for extraction).

\textbf{Orchestrators:}
\begin{itemize}
  \item \texttt{python experiments/run\_all.py} --- full extraction + core experiments; use \texttt{--skip-extraction} if \texttt{data/processed/} exists; \texttt{--figures-only} regenerates figures only; \texttt{--with-fewshot} adds long few-shot jobs.
  \item \texttt{python scripts/run\_value\_experiments.py --quick --figures --skip-sota --skip-exp1} --- ordered ``high-value'' pass for CPU (skips slow SOTA tuning and full Exp1; use \texttt{--skip-exp1} only when needed).
  \item \texttt{python experiments/exp1\_main.py --quick} --- fast Exp1 on CPU (sets single-thread BLAS before NumPy to avoid OpenMP deadlocks with PyTorch after sklearn/XGBoost). For the \emph{full} Exp1 protocol, run without \texttt{--quick} (GPU recommended).
\end{itemize}

\textbf{Outputs:} \texttt{python scripts/build\_thesis\_assets.py} regenerates figures, \texttt{ALL\_TABLES.tex}, \texttt{ALL\_TABLES\_APPENDIX.tex} (for Appendix~\ref{app:pipeline_tables}), and \texttt{KEY\_FINDINGS.md} (via \texttt{results/generate\_all\_outputs.py}). Before submission, run \texttt{python results/validate\_thesis\_readiness.py}.

Fix \texttt{random\_seed=42} in \texttt{config.yaml}. Hyperparameter grids are in Appendix~\ref{app:hyperparams}.
